\chapter{Lecture 21: Exception Safety, RAII and Smart Pointers}
\begin{center}
\textit{March 22, 2018}
\end{center}

\section{The Problem}
\begin{lstlisting}[style=cpp]
void f() {
    MyClass *p = new MyClass;
    MyClass mc;
    g();
    delete p;
}
\end{lstlisting}
Suppose \code{g} itself does not leak memory. Still, if \code{g} throws an
exception, \code{delete p} is never executed: \code{f} leaks memory.

We want to write \textbf{exception-safe} code. (This does \emph{not} mean ``code
that doesn't use exceptions''.) Exceptions are meant to be recovered from, so an
exception should not leave behind
\begin{itemize}
  \item memory leaks,
  \item dangling pointers,
  \item broken class invariants.
\end{itemize}

\begin{important}[C++ guarantee]
During stack unwinding, all stack-allocated objects are destroyed: their
destructors run and their memory is reclaimed. Nothing is promised about heap
memory.
\end{important}
So in \code{f}, \code{mc} is cleaned up, but the object \code{p} points to leaks.
A pointer is just a value: destroying \code{p} does not (and should not, on its
own) call \code{delete} on it.

\section{Refactoring by Hand}
\begin{lstlisting}[style=cpp]
void f() {
    MyClass *p = new MyClass;
    try {
        MyClass mc;
        g();
    } catch (...) {
        delete p;
        throw;          // re-throw: we only cleaned up
    }
    delete p;
}
\end{lstlisting}
This is tedious and error-prone (note the duplicated cleanup, and imagine five
pointers). We would like the ability to run some code no matter whether an
exception occurred:
\begin{itemize}
  \item Java has \code{finally};
  \item Scheme has \code{dynamic-wind}.
\end{itemize}
C++ has neither --- and doesn't need one, because of the guarantee above: the
only code guaranteed to run during unwinding is the destructors of stack
objects. So:
\begin{enumerate}
  \item maximize your use of the stack;
  \item but what if you really do need the heap?
\end{enumerate}

\section{RAII}
\begin{definition}
\textbf{RAII} --- \emph{Resource Acquisition Is Initialization} (a C++ idiom
named by Bjarne Stroustrup): every resource should be wrapped inside a
stack-allocated object whose destructor's job is to release the resource.
\end{definition}
We have been using it since the second C++ lecture:
\begin{lstlisting}[style=cpp]
{
    ifstream f{"hello.txt"};   // file is opened (acquired) at initialization
    ...
}                              // f's dtor runs: the file is closed (released)
\end{lstlisting}
The file is closed whether the block ends normally or by an exception.

Heap memory is a resource! Idea: wrap the heap object inside a stack object,
whose destructor deletes the heap object.

\section{\code{std::unique\_ptr}}
The standard library (\code{<memory>}) provides the template class
\code{std::unique\_ptr<T>}:
\begin{itemize}
  \item its ctor takes a \code{T*};
  \item its dtor calls \code{delete} on that pointer;
  \item it overloads \code{*} and \code{->}, so it can be used exactly like
        the raw pointer to the heap object.
\end{itemize}
\begin{lstlisting}[style=cpp]
#include <memory>

void f() {
    std::unique_ptr<MyClass> p{new MyClass};
    MyClass mc;
    g();
}   // no leak, whether or not g throws
\end{lstlisting}
Better: use \code{std::make\_unique}, which allocates the object and passes its
arguments to the \code{MyClass} ctor:
\begin{lstlisting}[style=cpp]
auto p = std::make_unique<MyClass>(/* args to MyClass ctor, if any */);
\end{lstlisting}

What about copying?
\begin{lstlisting}[style=cpp]
unique_ptr<C> p{new C};      // or: auto p = make_unique<C>();
unique_ptr<C> q = p;         // won't compile
\end{lstlisting}
If both \code{p} and \code{q} pointed to the object, it would be deleted twice.
So copying is disabled for \code{unique\_ptr}; moving still works (stealing is
still acceptable in C++ society):
\begin{lstlisting}[style=cpp]
unique_ptr<C> q = std::move(p);   // q owns the object; p is now null
\end{lstlisting}

\subsection{A Basic Implementation}
\begin{lstlisting}[style=cpp]
template <typename T> class unique_ptr {
    T *ptr;
  public:
    explicit unique_ptr(T *p) : ptr{p} {}
    ~unique_ptr() { delete ptr; }

    unique_ptr(const unique_ptr &other) = delete;              // no copy ctor
    unique_ptr &operator=(const unique_ptr &other) = delete;   // no copy assign

    unique_ptr(unique_ptr &&other) : ptr{other.ptr} {          // move ctor
        other.ptr = nullptr;
    }
    unique_ptr &operator=(unique_ptr &&other) {                // move assign
        using std::swap;
        swap(ptr, other.ptr);
        return *this;
    }

    T &operator*() const { return *ptr; }
    T *operator->() const { return ptr; }
    T *get() const { return ptr; }
};
\end{lstlisting}
\code{= delete} explicitly removes a function the compiler would otherwise
provide for free. (Note the move operations take \code{unique\_ptr\&\&}, not
\code{const unique\_ptr\&\&}: we must modify \code{other} to steal from it.)

\subsection{Ownership}
With smart pointers, the type of a pointer documents ownership:
\begin{itemize}
  \item \code{unique\_ptr<T>}: \textbf{owns} the object (responsible for deleting
        it). Composition (owns-a) relationships use \code{unique\_ptr} fields ---
        and then the class needs no hand-written destructor.
  \item raw pointer \code{T*}: \textbf{does not own}, just observes.
        Aggregation (has-a) relationships use raw pointers (or references). Get
        one from a \code{unique\_ptr} with \code{p.get()}.
  \item Never call \code{delete} on a raw pointer in a program that uses smart
        pointers for ownership.
\end{itemize}

\section{\code{std::shared\_ptr}}
Sometimes there is no single owner: several objects legitimately share one heap
object, and it should be deleted when the \emph{last} of them is done with it.
\begin{lstlisting}[style=cpp]
void f() {
    auto p1 = std::make_shared<C>();
    if (...) {
        auto p2 = p1;   // copy ctor: p1 and p2 point to the same heap object
        ...
    }                   // p2 goes out of scope: dtor runs, object NOT deleted
    ...
}                       // p1 goes out of scope: dtor runs, object deleted
\end{lstlisting}
Shared pointers use \textbf{reference counting}: a counter (stored on the heap
next to the object) records how many \code{shared\_ptr}s point to the same
object. Copying increments the count; destroying (or reassigning) a
\code{shared\_ptr} decrements it; when the count reaches 0 the object is deleted.

\begin{important}[Caveats]
\begin{itemize}
  \item Reference counting costs time and space; prefer \code{unique\_ptr} unless
        ownership really is shared.
  \item \textbf{Cycles}: if A holds a \code{shared\_ptr} to B and B holds a
        \code{shared\_ptr} to A, neither count ever reaches 0 and both leak.
        Break cycles by making one direction a non-owning pointer (raw pointer or
        \code{std::weak\_ptr}, which observes without contributing to the count).
  \item Never create two separate \code{shared\_ptr}s from the same raw pointer
        (\code{shared\_ptr<C> a\{p\}, b\{p\};}): they keep separate counts and
        both delete \code{p}. Use \code{make\_shared}, and copy the
        \code{shared\_ptr} itself.
\end{itemize}
\end{important}

\section*{Summary}
\begin{itemize}
  \item Exception safety: no leaks, dangling pointers or broken invariants when
        an exception passes through.
  \item C++ only guarantees that stack objects are destroyed during unwinding,
        so wrap every resource in a stack object (RAII).
  \item \code{unique\_ptr}: sole owner, move-only; create with
        \code{make\_unique}.
  \item \code{shared\_ptr}: shared ownership via reference counting; create with
        \code{make\_shared}; beware cycles.
  \item Raw pointers are for non-owning access.
\end{itemize}
